\section{Feasible-set universality and exact algebraic voltages}
\label{sec:algebraic}

The reduction preserves more than a yes-or-no answer. Its rational
solution-set correspondence transfers both topology and fields of definition.
We call two real sets \emph{rationally equivalent} if there is a
homeomorphism between them whose coordinate functions, and those of its
inverse, are rational functions with rational coefficients, with denominators
nonzero on their respective domains.

\begin{theorem}
\label{thm:universality}
Every compact semialgebraic set defined over $\Q$ is rationally equivalent
to the feasible voltage set of a rational $\RPF$ instance.
For every prescribed fixed girth bound, the instance can satisfy all the
structural and finite-data restrictions of Theorem~\ref{thm:structural}.
\end{theorem}
\begin{proof}
After clearing rational denominators, such a set is described by an ETR
formula with integer coefficients. Theorem 1 of
\citet{AbrahamsenMiltzow2019} gives a bounded $\ETRINV$ instance with a
rationally equivalent solution set. Compose that equivalence with
Proposition~\ref{prop:rational-bijection}, or with the stronger construction
of Theorem~\ref{thm:structural}. Composition preserves rationality and
continuity in both directions. Empty sets are included.
\end{proof}

This is an application of the arithmetic universality theorem, not a new
universality theorem for $\ETRINV$. It shows, for example, that restricting
to the graphs in Theorem~\ref{thm:structural} does not require the voltage
feasible set to be connected, contractible, or of dimension zero. Any compact
semialgebraic topology already occurs. The corresponding AC magnitude sets
inherit the same conclusion under the exact equal-angle transfers. For
rectangular voltages, fixing a reference angle in each connected component
removes the otherwise arbitrary common phase.

\begin{theorem}
\label{thm:algebraic-degree}
Let $\alpha$ be a real algebraic number. There is a rational $\RPF$
instance with a unique feasible voltage vector $v$ such that every $v_i$
belongs to $\Q(\alpha)$ and $\Q(v_j)=\Q(\alpha)$ for some designated bus
$j$. Consequently, a uniquely feasible rational network can have a designated
bus voltage of any prescribed algebraic degree over $\Q$.
All structural restrictions of Theorem~\ref{thm:structural} can be imposed.
The same statement holds for unique AC magnitudes, and for unique rectangular
voltages in the reference-fixed bus-angle-box model.
\end{theorem}
\begin{proof}
First suppose $\alpha\notin\Q$. Choose a nonzero polynomial in $\Q[t]$
having $\alpha$ as a root and rational endpoints $a<b$ whose interval
contains $\alpha$ and no other root. Its root condition together with
$a\le t\le b$ defines the singleton $\{\alpha\}$.
Theorem 1 of \citet{AbrahamsenMiltzow2019} gives a bounded $\ETRINV$
solution set rationally equivalent to this singleton, so it too is a
singleton, say $\{x\}$. Rationality of the forward map gives
$x_i\in\Q(\alpha)$ for every coordinate.

We also use the \emph{linear extension} conclusion of that theorem
(Definition 5 there): an original coordinate is recovered by a rational
affine rescaling of a designated replacement coordinate. Thus
$\alpha=A x_j+B$ for some $A,B\in\Q$. Since $\alpha$ is irrational,
$A\ne0$, and hence $\Q(x_j)=\Q(\alpha)$.
The network construction has a unique extension of $x$, with rational
coordinate functions, and it retains $x_j$ as the voltage at its original
variable root. The planarization, connectors, and subdivisions retain that
root coordinate and add only uniquely determined rational functions of the
old coordinates. This proves the claims for irrational $\alpha$.

For rational $\alpha$, a single voltage-$1$, zero-injection bus is sufficient:
its unique voltage generates $\Q=\Q(\alpha)$ and its graph meets every
stated restriction. To obtain degree $d\ge2$, take the positive real root
of $t^d-2$; this polynomial is irreducible by Eisenstein's criterion at $2$.
Degree one is the preceding rational example. Finally, the exact AC transfers
preserve precisely the RPF magnitude set, and reference-fixed angle boxes
force every phase to zero.
\end{proof}

The designated-coordinate assertion uses the affine recovery property,
not merely preservation of the field generated by all coordinates. The
theorem is an existence statement and asserts no additional size bound on
the affine recovery coefficients. It also does not by itself exclude
polynomially verifiable symbolic certificates: the relevant obstruction to
such certificates is the complexity consequence in
Corollary~\ref{cor:rpf-np}.
